\chapter{Sharp Runtime: Strings, Streams, and a Real Audit History}
\label{ch:sharp-runtime-namespaces}

This chapter covers two smaller, practical corners of \texttt{sharp-runtime}'s API --- its
text and stream types --- and then the part of its history most worth learning from: what
happened when the project's own completed task tracker was deliberately \emph{not} trusted as
proof of correctness, and independent audits went looking for what it might have missed.

\section{String and StringBuilder: a static-only utility versus a real mutable buffer}

\cnaclass{System::String} is deliberately \emph{not} a wrapper class at all --- its default
constructor and destructor are both explicitly deleted, making it a pure static-method utility
operating on ordinary \texttt{std::string} values (\texttt{IsNullOrEmpty}, \texttt{Compare},
the \texttt{IndexOf} / \texttt{LastIndexOf} family, \texttt{Split}, \texttt{Substring},
\texttt{Trim*}, \texttt{PadLeft} / \texttt{PadRight}). \cnaclass{Text::StringBuilder}, by
contrast, is a real, stateful, mutable buffer type --- the many-overloaded \texttt{Append}
(string, char, int, double, float, bool, and the \texttt{longcs} alias from
Chapter~\ref{ch:design-philosophy}), \texttt{AppendLine}, \texttt{Insert},
\texttt{Remove}, \texttt{Replace}, and a nested \texttt{ChunkEnumerator} class mirroring real
.NET's own chunked-buffer enumeration shape --- real .NET's own \texttt{StringBuilder} is
internally a linked list of separate character chunks, and \texttt{ChunkEnumerator} walks that
list without copying; this port's \texttt{StringBuilder} is backed by a single
\texttt{std::string} instead, so its own \texttt{ChunkEnumerator} always yields exactly one
chunk holding the entire current buffer --- the enumeration \emph{shape} is preserved
faithfully even though the simpler single-buffer implementation underneath never actually
produces more than one chunk to walk.

\begin{lstlisting}
// Worked example: System::String is a pure static utility -- there is no
// instance to construct, only free functions operating on std::string.
std::vector<std::string> parts = System::String::Split(
    "Content/Textures/hero.png", '/');
std::string trimmed = System::String::Trim(parts.back());

// Text::StringBuilder, by contrast, is a real mutable buffer.
System::Text::StringBuilder sb;
sb.Append("Score: ").Append(playerScore).AppendLine();
sb.Append("Lives: ").Append(livesRemaining);

for (const std::string& chunk : sb.GetChunks())
{
    // Exactly one iteration, always -- see the ChunkEnumerator note above.
    RenderHudText(chunk);
}
\end{lstlisting}

\section{Stream: a lightweight subset, honest about it}

\cnaclass{IO::Stream}'s own class comment states its scope plainly: ``lightweight subset of
the .NET Stream API.'' \texttt{Read(buffer, offset, count)}, \texttt{Close()}, and
\texttt{getLengthProperty()} are pure virtual; \texttt{Write} / \texttt{WriteByte} default to
throwing \texttt{NotSupportedException} unless a writable stream subclass overrides them ---
the same shape real .NET's own read-only stream implementations follow, rather than forcing
every stream subclass to implement a full read/write/seek contract regardless of whether it
actually supports all three. \texttt{getCanWriteProperty()} defaults to false and
\texttt{getCanReadProperty()} defaults to true, matching a read-only stream's honest self-report
without requiring every subclass to override both just to state the obvious.

\begin{lstlisting}
// Worked example: the minimal real subclass this base class actually
// requires -- only three overrides, matching a genuinely read-only
// stream's real capabilities (Write/WriteByte are correctly inherited as
// NotSupportedException-throwing, never called by a well-behaved reader).
class MemoryBlobStream final : public System::IO::Stream
{
public:
    explicit MemoryBlobStream(std::vector<SharpRuntime::bytecs> data)
        : data_(std::move(data)) {}

    SharpRuntime::intcs Read(SharpRuntime::bytecs buffer[],
                             SharpRuntime::intcs offset, SharpRuntime::intcs count) override
    {
        SharpRuntime::intcs n = std::min<SharpRuntime::intcs>(
            count, static_cast<SharpRuntime::intcs>(data_.size()) - position_);
        std::memcpy(buffer + offset, data_.data() + position_, n);
        position_ += n;
        return n;
    }

    void Close() override {}

    [[nodiscard]] SharpRuntime::intcs getLengthProperty() const override
    {
        return static_cast<SharpRuntime::intcs>(data_.size());
    }

private:
    std::vector<SharpRuntime::bytecs> data_;
    SharpRuntime::intcs position_ = 0;
};
\end{lstlisting}

\subsection{\cnaclass{MemoryStream}, and a small, deliberately-left divergence in its own constructor}

\cnaclass{System::IO::MemoryStream} is the concrete \texttt{Stream} subclass this ecosystem
actually uses in practice --- CNA's own \texttt{ContentManager} reaches for it to wrap data
loaded from sources like Android APK assets, where a real filesystem path is not available at
all. Its default constructor creates an empty, writable buffer, matching real .NET exactly. The
single-buffer constructor is where a small, real, currently accepted divergence lives:

\begin{lstlisting}
MemoryStream(const bytecs* buffer, intcs size);   // real signature, MemoryStream.hpp
\end{lstlisting}

Real .NET's own equivalent single-array constructor, \texttt{MemoryStream(byte[] buffer)},
defaults to a \emph{writable} stream --- its own documented implementation is literally
\texttt{this(buffer, true)}, forwarding to the two-argument overload with \texttt{writable}
explicitly \texttt{true}. This port's constructor instead defaults \texttt{writable\_} to
\texttt{false}, confirmed directly in \texttt{MemoryStream.cpp}:

\begin{lstlisting}
MemoryStream::MemoryStream(const bytecs* buffer, intcs size)
    : data_(buffer, buffer + size), position_(0), writable_(false) {}
\end{lstlisting}

This was found, not invented, incidentally while writing regression tests for an unrelated
audit ticket --- and deliberately left as-is right there rather than fixed on the spot, to
avoid scope creep into a ticket that was not about this constructor at all. It remains a
real, open, precisely-scoped divergence today: code that passes a raw buffer to this
constructor and then calls \texttt{Write} on the result will get a real
\texttt{NotSupportedException} in this port, where the equivalent real .NET call would have
succeeded. A caller that genuinely wants a writable copy of existing buffer contents should
construct an empty \cnaclass{MemoryStream} and \texttt{Write} the source bytes into it
explicitly, rather than relying on this specific constructor's own writability default.

\subsection{Closing a memory stream closes operations, not the bytes it already owns}

The constructor difference above should not be confused with this class's much more carefully
matched close/dispose behavior.  \texttt{MemoryStream::Close()} changes only its
\texttt{isOpen\_} state; it deliberately retains both the vector of bytes and the current
position.  After close, \texttt{Read}, \texttt{Write}, \texttt{WriteByte}, \texttt{Seek},
\texttt{Length}, and \texttt{Position} correctly reject use with
\cnaclass{ObjectDisposedException}, while \texttt{CanSeek} returns false.  In deliberate
contrast, \texttt{ToArray()} and \texttt{GetBuffer()} remain usable:

\begin{lstlisting}
System::IO::MemoryStream stream;
const SharpRuntime::bytecs payload[] = {1, 2, 3, 4};
stream.Write(payload, 0, 4);
stream.Close();

std::vector<SharpRuntime::bytecs> saved = stream.ToArray();  // still {1,2,3,4}
// stream.getLengthProperty();  // throws ObjectDisposedException
// stream.Read(buffer, 0, 1);   // throws ObjectDisposedException
\end{lstlisting}

That asymmetry is intentional, not a use-after-dispose loophole.  Real .NET's own
\texttt{MemoryStream.Dispose(bool)} explicitly preserves its buffer so post-disposal
\texttt{TryGetBuffer}, \texttt{GetBuffer}, and \texttt{ToArray} can still retrieve data; the
runtime's \texttt{ensureNotClosed()} guard follows the same rule by guarding operations but not
the two extraction methods.  CNA callers can therefore finish producing an in-memory content
blob, close the stream to catch accidental later I/O, and still hand a stable copy to a consumer.
The distinction between the methods matters: \texttt{ToArray()} returns an independent vector,
safe across later stream writes, whereas \texttt{GetBuffer()} is a reference to the live vector
and may be invalidated by a reallocation.

The same implementation also has two less obvious safety properties whose tests prevent a
``memory only'' stream from becoming an unchecked byte container.  Null buffers and negative
offsets/counts are rejected with the real \cnaclass{ArgumentNullException} or
\cnaclass{ArgumentOutOfRangeException} rather than being mistaken for an EOF return of zero.
And a caller may legally set \texttt{Position} past the current end, but a later write computes
\texttt{position + count} in \texttt{int64\_t} before resizing.  If that sum exceeds the
runtime's signed \texttt{intcs} range it throws \cnaclass{IO::IOException} (``Stream was too
long.'') instead of wrapping negative and writing beyond the vector.  This was not defensive
theorizing: the former narrow-integer calculation had the same signed-overflow shape already
reproduced under UBSan in a related slicing path, and \texttt{StreamTests.cpp} now exercises
\texttt{Position = 2147483647} followed by a ten-byte write specifically to require the
exception rather than an allocation or memory corruption.

\section{\cnaclass{UTF8Encoding}: real conformance validation, an untested fallback path}

\cnaclass{Text::UTF8Encoding}'s own class comment states a real, specific guarantee that is
easy to assume any UTF-8 encoder provides and that this port genuinely earns rather than
skips: \texttt{GetBytes()} / \texttt{GetString()} \emph{validate} that input is well-formed
UTF-8 --- not merely pass bytes through as opaque \texttt{char} values --- and route any
ill-formed byte at each invalid position through the configured fallback before resuming.
Reading \texttt{UTF8Encoding.cpp}'s own local \texttt{wellFormedUtf8Length()} helper directly
shows this is real conformance checking, not a length-only byte count: a leading byte is
walked against the real UTF-8 grammar (1/2/3/4-byte forms, each continuation byte checked for
the \texttt{10xxxxxx} pattern), and even a \emph{structurally} well-formed sequence is rejected
if the code point it decodes to is an overlong encoding (a multi-byte sequence encoding a code
point that a shorter form could already represent --- a real, historically security-relevant
UTF-8 attack shape, since a naive decoder that accepts overlong encodings lets an attacker
smuggle a byte value like \texttt{0x2F} (\texttt{'/'}) past a filter checking only the raw byte
stream) or a UTF-16 surrogate code point (\texttt{U+D800}--\texttt{U+DFFF}, which UTF-8 must
never directly encode, since surrogates exist only as a UTF-16 encoding artifact).

\begin{lstlisting}
// From wellFormedUtf8Length() -- the 3-byte case, real conformance logic,
// not just "are there two continuation bytes after this lead byte":
if ((c0 & 0xF0) == 0xE0 && /* two real continuation bytes follow */) {
    uint32_t cp = /* decode the 16-bit code point from all three bytes */;
    return (cp >= 0x800 && !(cp >= 0xD800 && cp <= 0xDFFF)) ? 3 : 0;
    //      ^ rejects overlong (cp < 0x800 has a shorter valid encoding)
    //                          ^ rejects UTF-16 surrogates directly encoded in UTF-8
}
\end{lstlisting}

The constructor also earns a real, specific correction over the naive default: real .NET's own
\texttt{UTF8Encoding.SetDefaultFallbacks()} substitutes the Unicode replacement character
\texttt{U+FFFD} for invalid input, not the generic \cnaclass{Encoding} base class's plain
\texttt{"?"} default (correct only for single-byte code pages, where every valid character
already has an unambiguous byte representation to fall back to). This port's constructor
reproduces that distinction precisely:

\begin{lstlisting}
UTF8Encoding::UTF8Encoding() {
    // Real UTF8Encoding.SetDefaultFallbacks() uses a U+FFFD replacement
    // fallback, not the generic Encoding base class's "?" default.
    setEncoderFallbackProperty(
        std::make_shared<EncoderReplacementFallback>("\xEF\xBF\xBD"));
    setDecoderFallbackProperty(
        std::make_shared<DecoderReplacementFallback>("\xEF\xBF\xBD"));
}
\end{lstlisting}

A worked example makes the whole mechanism concrete --- three real byte sequences, one
well-formed and two deliberately malformed in different ways, through the same
\texttt{GetString()} call:

\begin{lstlisting}
System::Text::UTF8Encoding utf8;

// Well-formed: a real 3-byte sequence for U+20AC (EURO SIGN).
const SharpRuntime::bytecs euroSign[] = {0xE2, 0x82, 0xAC};
std::string ok = utf8.GetString(euroSign, 0, 3);          // EURO SIGN, unchanged

// Overlong encoding of U+002F ('/') using the 2-byte form -- structurally
// valid continuation-byte shape, but wellFormedUtf8Length() rejects it
// because cp (0x2F) is below the 2-byte form's own 0x80 minimum.
const SharpRuntime::bytecs overlongSlash[] = {0xC0, 0xAF};
std::string rejected = utf8.GetString(overlongSlash, 0, 2);  // two U+FFFD

// A lone continuation byte with no lead byte before it -- fails every one
// of wellFormedUtf8Length()'s four branches, falls through to the final
// "return 0" unconditionally.
const SharpRuntime::bytecs loneContinuation[] = {0x80};
std::string alsoRejected = utf8.GetString(loneContinuation, 0, 1);  // one U+FFFD
\end{lstlisting}

The overlong and lone-continuation cases each substitute one \texttt{U+FFFD} per invalid
\emph{byte}, not per attempted sequence --- \texttt{overlongSlash}'s two bytes both fail
validation independently (the loop advances one byte at a time through the fallback path, per
\texttt{UTF8Encoding.cpp}'s own \texttt{++i} in the non-well-formed branch), so a 2-byte
overlong sequence becomes two replacement characters, not one. This matches real .NET's own
observable behavior for the identical input, confirmed by working through the .NET runtime's
own documented per-byte fallback-and-resynchronize algorithm for UTF-8 decoding rather than
assumed by analogy.

None of this specific malformed-input behavior has a dedicated regression test anywhere in the
project today --- \texttt{EncodingTests.cpp}'s own \texttt{UTF8*} cases cover only well-formed
round-trips (\texttt{UTF8GetBytesHello}, \texttt{UTF8RoundTrip}, empty-string and offset
handling), never an overlong encoding, a bare continuation byte, a truncated multi-byte
sequence, or a directly-encoded surrogate. The conformance logic above is real and correctly
grounded in the UTF-8 grammar by direct code reading, but --- consistent with this chapter's
own recurring theme --- its correctness currently rests entirely on that reading, not on a
passing test asserting it.

\section{The post-stabilization audit: distrusting a complete-looking task list on purpose}

\texttt{sharp-runtime}'s own ticket backlog --- 1{,}709 tickets across sixteen categories ---
reached 100\% \texttt{done} at one point in the project's history. What happened next is the
single most instructive episode in this project's own story, and worth recounting in full,
because it is the same discipline Chapter~\ref{ch:project-practice} named as this whole
ecosystem's most distinctive trait, applied at real scale: rather than treating a fully
checked-off ticket table as proof the runtime was correct, the project commissioned a fresh,
independent audit specifically built to distrust that conclusion. Seven parallel, read-only
audit agents were dispatched across eight risk categories --- API inconsistencies,
undocumented deviations from real .NET, remaining \texttt{NotImplementedException}s, partial
or stub classes, platform-specific issues, exception-type mismatches, silently-wrong
behavior, and missing high-risk test coverage --- each one required to verify every finding
directly against real .NET source and against the project's own documented, accepted
deviations list before it could be reported at all, specifically to filter out
plausible-sounding but wrong findings before they ever reached a human.

The result was twenty concrete findings --- nine high-severity, eight medium, three low ---
filed as new tickets and subsequently fixed. One is worth naming in full, because of how
quietly it could have shipped: \texttt{Convert::ToXxx(double/float)} was implemented as a
bare \texttt{static\_cast<T>()}, meaning simple truncation rather than round-half-to-even
rounding. A direct repro confirmed \texttt{Convert::ToInt32(2.9)} returned \texttt{2} where
real .NET returns \texttt{3}, and \texttt{Convert::ToInt32(3.5)} returned \texttt{3} where real
.NET returns \texttt{4} --- and, more strikingly, \texttt{ToUInt32(double)}'s own doc comment
already said \texttt{"(truncates)"}, as though the wrong behavior had been documented as
intentional rather than caught as a bug.

\begin{lstlisting}
// Worked example: the exact repro that caught this bug, and what
// Convert::ToInt32(double) returns now that it implements real
// round-half-to-even rounding instead of a bare static_cast<int>().
assert(System::Convert::ToInt32(2.9) == 3);  // was 2 before the fix
assert(System::Convert::ToInt32(3.5) == 4);  // was 3 before the fix -- .NET
                                              // rounds 3.5 up to 4, not to
                                              // the nearer *even* integer,
                                              // since 4 is already even
\end{lstlisting}

\section{Two more findings from the same audit, worth their own worked examples}

Two further findings from that first audit round are worth their own concrete treatment,
because each illustrates a distinct failure shape beyond simple truncation, and this project's
own source shows both bugs and both fixes side by side today.

\cnaclass{Dictionary<TKey, TValue>}'s non-\texttt{const} indexer used to be a plain
\texttt{TValue\&}, which could not distinguish a read access from a write access --- so reading
a missing key silently inserted a default-constructed value, matching
\texttt{std::unordered\_map::operator[]}'s own convention rather than real .NET's, where the
indexer getter unconditionally throws \texttt{KeyNotFoundException} on any missing key. The
quietly damning detail the audit surfaced: this exact bug had already been found and fixed once
in this same codebase, in \cnaclass{Concurrent::ConcurrentDictionary}, via a dedicated
\texttt{ValueProxy} wrapper whose own doc comment states plainly that it does \emph{not} insert
a default --- the fix had simply never been propagated to the plain, non-concurrent
\cnaclass{Dictionary}. It has been now:

\begin{lstlisting}
// Worked example: Dictionary's operator[] today, a ValueProxy return type
// (not a plain TValue&) so a read and a write can be told apart, exactly
// mirroring the precedent ConcurrentDictionary::ValueProxy already set.
[[nodiscard]] ValueProxy operator[](const TKey& key) { return ValueProxy(this, key); }

[[nodiscard]] const TValue& operator[](const TKey& key) const
{
    auto it = map_.find(key);
    if (it == map_.end())
        throw KeyNotFoundException("The given key was not present in the dictionary.");
    return it->second;
}
\end{lstlisting}

\cnaclass{Concurrent::ConcurrentDictionary::GetOrAdd} / \texttt{AddOrUpdate} had a second, more
subtle problem the audit also caught: both held their internal \texttt{std::mutex} for the
\emph{entire} call, including the user-supplied factory callback's own invocation. Since
\texttt{std::mutex} is non-recursive, a factory that reentrantly called back into the very same
dictionary instance --- a realistic memoization pattern, not a contrived edge case --- would
deadlock the calling thread against itself. The fix moves the factory call outside the lock
entirely, matching real .NET's own documented contract for this method precisely (the factory
may run without the lock held and, under real contention, may even be invoked more than once,
with the loser's result discarded):

\begin{lstlisting}
// Worked example: GetOrAdd today -- the factory call is genuinely outside
// both lock_guard scopes, so a reentrant factory can no longer self-deadlock.
TValue GetOrAdd(const TKey& key, std::function<TValue(const TKey&)> factory)
{
    {
        std::lock_guard<std::mutex> lk(mutex_);
        auto it = map_.find(key);
        if (it != map_.end()) return it->second;
    }
    TValue v = factory(key);   // deliberately unlocked -- see the note above
    std::lock_guard<std::mutex> lk(mutex_);
    auto [it, inserted] = map_.emplace(key, std::move(v));
    return it->second;
}
\end{lstlisting}

\section{A systemic finding revisited: fail-fast enumeration, fixed in ten of eleven collections}

One finding from the same first audit round deserves its own treatment separately from the two
above, because of its scope: \texttt{POST\_STABILIZATION\_AUDIT.md} flagged that \emph{every}
\cnans{System.Collections.Generic} mutable collection --- eleven files,
\cnaclass{List}\allowbreak{}\texttt{<T>} through \cnaclass{PriorityQueue} --- lacked the
version-counter fail-fast check real .NET's own generic collections all share: mutate a
collection while an enumerator from it is still live, and the very next \texttt{MoveNext()}
should throw \texttt{InvalidOperationException} (``Collection was modified; enumeration
operation may not execute'') rather than silently skipping elements or, worse, dereferencing a
dangling iterator into a reallocated buffer. The audit noted the fix pattern already existed
elsewhere in this same codebase (the legacy, non-generic \texttt{Hashtable} / \texttt{ArrayList}
/ \texttt{Queue} / \texttt{Stack} types all had it) --- it had simply never been extended to the
far more heavily used generic collections, and its own scope note predicted a real fix would
need per-type application across all eleven files.

Reading the current headers directly rather than trusting either the audit's own snapshot or a
completion claim shows the fix landed, and landed almost completely: \cnaclass{List}\allowbreak{}
\texttt{<T>}'s own \texttt{Enumerator} now captures \texttt{version\_} at construction and checks
it on every \texttt{MoveNext()} / \texttt{Reset()}, throwing exactly the real .NET message:

\begin{lstlisting}
class Enumerator : public IEnumerator<T>
{
    const List<T>* list_;
    intcs version_;
    intcs index_ = -1;
public:
    explicit Enumerator(const List<T>* list) : list_(list), version_(list->version_) {}
    bool MoveNext() override {
        if (version_ != list_->version_)
            throw System::InvalidOperationException(
                "Collection was modified; enumeration operation may not execute.");
        return ++index_ < static_cast<intcs>(list_->items_.size());
    }
    // ...
};
\end{lstlisting}

Grepping every one of the audit's own eleven named files for a real \texttt{version\_} field
confirms the same pattern is now present in ten of them ---
\cnaclass{List}\allowbreak{}\texttt{<T>}, \cnaclass{Dictionary}\allowbreak{}\texttt{<K,V>},
\cnaclass{HashSet}\allowbreak{}\texttt{<T>}, \cnaclass{LinkedList}\allowbreak{}\texttt{<T>},
\cnaclass{Queue}\allowbreak{}\texttt{<T>}, \cnaclass{Stack}\allowbreak{}\texttt{<T>},
\cnaclass{SortedDictionary}\allowbreak{}\texttt{<K,V>},
\cnaclass{SortedList}\allowbreak{}\texttt{<K,V>}, \cnaclass{SortedSet}\allowbreak{}\texttt{<T>},
and \cnaclass{OrderedDictionary} --- not a partial or cosmetic fix, a real, independently
verified per-type application matching exactly what the audit's own scope note anticipated
would be needed.

\cnaclass{List}\allowbreak{}\texttt{<T>}'s own class doc-comment is honest about one narrow gap
the fix does \emph{not} close, worth quoting directly rather than glossing over: real .NET's
index-assignment setter (\texttt{list[i] = value}) also bumps its internal version, but this
port's \texttt{operator[]} returns a plain \texttt{T\&} for C++ ergonomics, so there is no hook
point to intercept a later \texttt{list[i] = value;} assignment made through that reference ---
value-only index writes do not trigger fail-fast detection here. The comment explicitly weighs
and rejects the fix (widening \texttt{operator[]} into a proxy-object return, the same technique
already used to fix \cnaclass{Dictionary}\allowbreak{}\texttt{<K,V>}'s own indexer bug earlier in
this chapter) as too large and risky a change to the single most call-site-heavy method in the
codebase to make unilaterally --- a deliberately deferred, not overlooked, gap.

\cnaclass{PriorityQueue}\allowbreak{}\texttt{<TElement,TPriority>} is the one file of the eleven
with zero \texttt{version\_} occurrences, and it is worth checking \emph{why} before assuming the
fix was simply skipped there too. Reading the header directly shows a different, more fundamental
reason: this port's \cnaclass{PriorityQueue} does not implement any enumerable interface at all
--- no \texttt{GetEnumerator()}, no \texttt{begin()} / \texttt{end()}, and no equivalent of real
.NET's own \texttt{UnorderedItems} property, which is exactly what \texttt{PriorityQueue<TElement,
TPriority>} in real .NET exposes for enumeration. There is no enumerator here for a mutation to
invalidate, so the fail-fast pattern genuinely does not apply --- not a missed eleventh
application of finding \#4, but a separate, real gap of its own (no enumeration support at all),
one this chapter's own source-grounding discipline requires stating precisely rather than folding
into the version-tracking story just because both would show up in the same file-by-file grep.

Two smaller items from the identical audit round are worth a one-line confirmation each, both
independently re-verified against current source: \cnaclass{MemoryStream}\allowbreak{}
\texttt{::Write} previously validated only \texttt{buffer == nullptr} and \texttt{count <= 0},
letting a negative \texttt{offset} reach an unchecked \texttt{std::copy} and read one byte before
the caller's own array --- confirmed via a real ASan repro, per the fix's own code comment, and
now closed with explicit \texttt{ArgumentNullException} / \texttt{ArgumentOutOfRangeException}
checks matching \texttt{Read()}'s own validation in the same class. And
\cnaclass{List}\allowbreak{}\texttt{<T>}'s doc-comment claim of ``full \texttt{IList<T>}
compliance'' while \texttt{CopyTo} was entirely missing --- also fixed, with the exact two
overloads real .NET's own \texttt{ICollection<T>} contract specifies.

\section[A second audit, and what sanitizers found that reasoning alone did not]{A second, independent audit, and what sanitizers found that reasoning alone did not}

A further, separate audit produced seven more findings, six fixed in the same session ---
among them a stream-backed ZIP archive implementation that never actually wrote its output
back to the underlying stream in create or update mode, and five of this project's own
destructors (spanning compression streams and XML writing) that could call
\texttt{std::terminate()} if an I/O failure occurred during RAII cleanup, since a
C\texttt{++} destructor that lets an exception escape terminates the program outright rather
than propagating the failure the way a .NET \texttt{Dispose()} call safely could.

A third of that same round's findings deserves its own worked treatment, because it is a real
race between two things that individually look correct: \cnaclass{TaskCompletionSource<TResult>}%
\texttt{::TrySetResult} first atomically claims completion (\texttt{completed\_.%
compare\_exchange\_strong}), \emph{then} calls \texttt{promise\_.set\_value(result)} to actually
publish the result. If \texttt{TResult}'s own copy constructor throws --- a real possibility for
any non-trivial result type, not a contrived one --- that throw happens \emph{after} completion
was already claimed but \emph{before} the underlying \texttt{std::promise}'s shared state ever
becomes ready. Before the fix, nothing caught that exception at this layer, so it propagated
straight out of \texttt{TrySetResult} with the promise left permanently unset: every current and
future waiter on \texttt{GetResult()} or the bridging \cnaclass{Task} blocked forever, and no
later \texttt{TrySetResult} / \texttt{TrySetException} / \texttt{TrySetCanceled} call could ever
re-claim \texttt{completed\_} to unblock them, since it was already (falsely) true.

\begin{lstlisting}
bool TrySetResult(const TResult& result) {
    bool expected = false;
    if (!completed_.compare_exchange_strong(expected, true)) return false;
    try {
        promise_.set_value(result);         // TResult's copy ctor could throw here
    } catch (...) {
        promise_.set_exception(std::current_exception());  // settle the promise anyway
        throw;                              // still propagates to this immediate caller
    }
    return true;
}
\end{lstlisting}

The fix settles the promise with the copy failure itself inside the same \texttt{catch}, so
every waiter still unblocks --- now with an exception to observe rather than a hang --- while
the immediate caller of \texttt{TrySetResult} still sees the same throw it always did, keeping
existing callers that already wrap the call in their own \texttt{try} / \texttt{catch} unaffected.
The general shape worth remembering: an atomic ``claim the transition'' step and the actual
state-publishing step that follows it are two separate operations, and a real exception thrown
strictly between them is exactly the kind of failure window unit tests built around the happy
path will never exercise on their own.

A first-ever full run of the full sanitizer trio --- ThreadSanitizer, AddressSanitizer, and
UndefinedBehaviorSanitizer --- found real bugs no amount of code review alone had caught:
ThreadSanitizer found one genuine production deadlock (a \texttt{Channel} implementation
missing a \texttt{notify\_all} call, a lost-wakeup bug affecting multiple blocked waiters)
alongside three test-only races; AddressSanitizer found one confirmed heap-buffer-overflow (an
aligned-reallocation helper copying the \emph{old}, smaller allocation's size into a newly
enlarged buffer) plus five real memory leaks in XML document node creation; and
UndefinedBehaviorSanitizer found one genuine undefined-behavior call reaching into vendored
\texttt{miniz} code, from an empty vector's \texttt{.data()} pointer reaching a
non-null-parameter-annotated function.

\section{Why this history belongs in a book, not just a changelog}

None of these bugs are individually remarkable --- a truncation-instead-of-rounding bug, a
missed \texttt{notify\_all}, a leaking XML method. What is worth taking from this chapter is
the project's own response to having a fully completed task tracker: rather than reading
``100\% done'' as ``verified correct,'' it treated that state as precisely the moment a fresh,
independently-skeptical audit was most valuable, and built that audit specifically to check
claims against real .NET source rather than against the project's own prior conclusions. This
is the same discipline \texttt{xna4-spec} auditing (Chapter~\ref{ch:xna4-spec-auditing})
brings to CNA's own XNA surface, applied here one layer further down, to the foundation layer
everything else in this book ultimately stands on.
